% 火星（综述）
% license CCBYSA3
% type Wiki

本文根据 CC-BY-SA 协议转载翻译自维基百科\href{https://en.wikipedia.org/wiki/Mars}{相关文章}。

\begin{figure}[ht]
\centering
\includegraphics[width=6cm]{./figures/b86f152bd3571454.png}
\caption{火星的真实颜色\(^{\text{a}}\)，由“希望号”轨道器拍摄。画面中央可见塔尔西斯山脉，左侧为奥林匹斯山，右侧则是水手谷。} \label{fig_Mars_1}
\end{figure}
火星是距离太阳的第四颗行星。由于其呈橙红色的外观，它也被称为“红色星球”\(^{\text{[22,23]}}\)。火星是一颗类似沙漠的岩质行星，大气稀薄，主要成分为二氧化碳（CO\(_2\)）。在平均地表高度处，其大气压仅为地球的几千分之一，大气温度范围为 −153 至 20 °C（−243 至 68 °F）\(^{\text{[24]}}\)，宇宙辐射水平很高。火星保留了一部分水，存在于地下以及稀薄大气中，可形成卷云、雾气、霜冻、较大的极地永久冻土区与冰帽（伴随季节性的 CO\(_2\) 雪），但其表面不存在液态水体。其表面重力约为地球的三分之一，为月球的两倍。直径为 6,779 公里（4,212 英里），约为地球的一半，是月球的两倍，其表面积则相当于地球全部陆地面积的总和。

细微尘埃遍布火星表面和大气，在低重力条件下，即便是稀薄大气的微弱风力也能将其扬起并扩散开来。火星地形大致呈现显著的南北分界，即“火星二元性结构”：北半球主要是平坦、低洼的平原，而南半球则由遍布陨石坑的高地组成。从地质学上看，火星处于相对活跃的状态，地下会发生“火震”，但其上也分布着许多巨型而已灭绝的火山（最高的是奥林匹斯山，高 21.9 公里或 13.6 英里），以及太阳系中最大峡谷之一的水手谷（长约 4,000 公里或 2,500 英里）。火星有两颗天然卫星，形状小而不规则：火卫一（Phobos）与火卫二（Deimos）。由于其轴倾角为 25 度，火星像地球（轴倾角 23.5 度）一样经历季节变化。一个火星太阳年相当于 1.88 个地球年（687 个地球日），一个火星太阳日（sol）则等于 24.6 小时。

火星与其他行星一同形成于大约 45 亿年前。在火星的诺亚纪（Noachian period，约 45 亿年至 35 亿年前）期间，其地表经历了陨石撞击、峡谷形成、侵蚀作用、可能存在的海洋，以及磁层的消失。赫斯珀里亚纪（Hesperian period，自约 35 亿年前开始，结束于约 33–29 亿年前）以广泛的火山活动和导致巨大泄流河道形成的洪水为主要特征。亚马逊纪（Amazonian period）从其后延续至今，是当前地质过程的主要影响因素。由于火星的地质历史，对于火星过去或当前是否存在生命的探索仍是活跃的科学研究领域，其中一些可能的生命迹象仍需进一步确认。

由于在地球夜空中可被肉眼看到、呈红色并缓慢移动，火星自古以来就被人类观测，并在不同文化中具有多样的象征意义。1963 年，人类首次发射前往火星的探测器“火星 1 号”（Mars 1），但途中失联。第一项成功飞越火星的探测任务发生在 1965 年，由“水手 4 号”（Mariner 4）完成。1971 年，“水手 9 号”（Mariner 9）进入火星轨道，成为人类历史上首个进入除月球、太阳或地球以外天体轨道的航天器；同年，“火星 2 号”成为首个在火星发生非受控撞击的探测器，“火星 3 号”则实现了首次成功着陆。自 1997 年以来，探测器在火星上持续运行。有时，超过十个探测器会同时在火星轨道或地表开展活动，这一数量超过地球以外任意行星。

火星常被视为未来载人探索任务的目标，尽管目前尚无正式计划付诸实施。
\subsection{自然历史}  
\subsubsection{形成} 
科学家推测，在太阳系形成过程中，火星是由于围绕太阳运行的原行星盘中的物质通过随机的失控聚积过程形成的。火星在太阳系中的位置使其具备许多独特的化学特征。相较于地球，沸点较低的元素，例如氯、磷和硫，在火星上要常见得多；这些元素很可能是在年轻太阳强烈的太阳风作用下被推移到更外侧区域\(^{\text{[25]}}\)。

晚期重轰炸  
在行星形成之后，内太阳系可能经历了所谓的“晚期重轰炸”。火星约有 60\% 的表面保留了那个时期的撞击记录\(^{\text{[26,27,28]}}\)，而其余许多区域的地表之下，可能也存在由同一时期事件形成的巨大撞击盆地。然而，最近的建模研究对晚期重轰炸的存在提出了质疑\(^{\text{[29]}}\)。证据显示，火星北半球可能存在一个巨大的撞击盆地，范围达 10,600 × 8,500 公里（6,600 × 5,300 英里），约为月球南极—艾特肯盆地面积的四倍；若得到确认，这将是迄今发现的最大撞击盆地\(^{\text{[30]}}\)。科学家推测该盆地形成于约 40 亿年前，当时火星被一个与冥王星大小相当的天体撞击，造成了火星南北半球的显著地形差异，并形成覆盖火星 40\% 表面的平滑的北方平原（Borealis basin）\(^{\text{[31,32]}}\)。

一项 2023 年的研究根据火卫二（Deimos）的轨道倾角提出证据，显示火星在 35 亿至 40 亿年前可能曾拥有一个环系统\(^{\text{[33]}}\)。该环系统可能由一颗质量约为火卫一（Phobos）20 倍的古老卫星解体形成；而火卫一可能正是这一原始环系统的残余物\(^{\text{[34,35]}}\)。
\begin{figure}[ht]
\centering
\includegraphics[width=14.25cm]{./figures/4b6ac302fa80c765.png}
\caption{} \label{fig_Mars_2}
\end{figure}
\begin{figure}[ht]
\centering
\includegraphics[width=14.25cm]{./figures/7861217b62cae434.png}
\caption{} \label{fig_Mars_3}
\end{figure}
火星的地质历史可以被划分为多个时期，但以下三个时期最为主要 \(^{\text{[36,37]}}\)。
\begin{itemize}
\item 诺亚纪（Noachian period）：约 45 亿年至 35 亿年前，是火星现存最古老地表的形成阶段。诺亚纪地表布满大量大型撞击坑。塔尔西斯隆起（Tharsis bulge）这一火山高地被认为形成于该时期，其末期可能经历了大规模液态水泛滥。该时期名称源自 Noachis Terra \(^{\text{[38]}}\)。
\item 赫斯珀纪（Hesperian period）：约 35 亿年至介于 33 亿年至 29 亿年前之间。赫斯珀纪的特征是广泛的熔岩平原的形成。名称源自 Hesperia Planum \(^{\text{[38]}}\)。
\item 亚马孙纪（Amazonian period）：约 33 亿年至 29 亿年前至今。亚马孙纪区域的陨石撞击坑数量较少，但地貌表现多样。奥林匹斯山（Olympus Mons）形成于该时期，同时火星其他地区也存在熔岩流活动。名称源自 Amazonis Planitia \(^{\text{[38]}}\)。
\end{itemize}
\subsubsection{近期地质活动}
火星上仍然存在地质活动。阿萨巴斯卡峡谷（Athabasca Valles）中有形成于约 2 亿年前的片状熔岩流。刻耳柏洛斯裂谷（Cerberus Fossae）这一断陷带中的水流活动发生于不足 2,000 万年前，显示出同样年轻的火山侵入事件 \(^{\text{[39]}}\)。火星勘测轨道器（Mars Reconnaissance Orbiter）已拍摄到雪崩的影像 \(^{\text{[40,41]}}\)。
\subsection{物理特征}
\begin{figure}[ht]
\centering
\includegraphics[width=8cm]{./figures/23fb842f59355847.png}
\caption{按比例绘制的火星及内太阳系中具有行星质量的天体。从左至右依次为：水星、金星、地球、月球、火星和谷神星。} \label{fig_Mars_4}
\end{figure}
火星的直径约为地球的一半，或约为月球的两倍，其表面积仅略小于地球全部陆地面积 \(^{\text{[2]}}\)。火星的密度低于地球，体积约为地球的 15\%，质量约为地球的 11\%，因此其表面重力约为地球的 38\%。火星是目前已知的唯一一颗沙漠行星，即表面类似于地球荒漠环境的类地岩质行星。火星表面呈红橙色，是由于铁(III)氧化物（纳米级 Fe\(_2\)O\(_3\)）以及铁(III)氧化物-氢氧化物矿物针铁矿（goethite）所致 \(^{\text{[42]}}\)。其颜色有时类似奶油太妃糖 \(^{\text{[43]}}\)；其他常见的地表颜色还包括金色、棕色、褐色以及因矿物组成不同而呈现的绿ish色调 \(^{\text{[43]}}\)。
\subsubsection{内部结构}
\begin{figure}[ht]
\centering
\includegraphics[width=8cm]{./figures/4694e492a3c0d60d.png}
\caption{火星的内部结构 \(^{\text{[44,45]}}\)} \label{fig_Mars_5}
\end{figure}



